
\documentclass{article}
\usepackage{colortbl}
\usepackage{makecell}
\usepackage{multirow}
\usepackage{supertabular}

\begin{document}

\newcounter{utterance}

\centering \large Interaction Transcript for game `bbh', experiment `salient\_translation\_error\_detection', episode 46 with Qwen3.6{-}27B.
\vspace{24pt}

{ \footnotesize  \setcounter{utterance}{1}
\setlength{\tabcolsep}{0pt}
\begin{supertabular}{c@{$\;$}|p{.15\linewidth}@{}p{.15\linewidth}p{.15\linewidth}p{.15\linewidth}p{.15\linewidth}p{.15\linewidth}}
    \# & \multicolumn{2}{c}{Player} && \multicolumn{2}{c}{Game Master} \\
    \hline

    \theutterance \stepcounter{utterance}  
    & & & \multicolumn{4}{p{0.6\linewidth}}{
        \cellcolor[rgb]{0.9,0.9,0.9}{
            \makecell[{{p{\linewidth}}}]{
                \texttt{\tiny{[P1$\langle$GM]}}
                \texttt{Detect the type of error in an English translation of a German source sentence.} \\
\\ 
\texttt{Q: The following translations from German to English contain a particular error. That error will be one of the following types: Named Entities: An entity (names, places, locations, etc.) is changed to a different entity. Numerical Values: Numerical values (ordinals or cardinals), dates, and/or units are changed. Modifiers or Adjectives: The modifiers and adjectives pertaining to a noun are changed. Negation or Antonyms: Introduce or remove a negation or change comparatives to their antonyms. Facts: Trivial factual errors not pertaining to the above classes are introduced in the translations. Dropped Content: A significant clause in the translation is removed. Please identify that error.  Source: In der Liste der Baudenkmale in Lenzen (Elbe) sind alle Baudenkmale der brandenburgischen Stadt Lenzen (Elbe) und ihrer Ortsteile aufgelistet.} \\
\texttt{Translation: In the list of architectural monuments in Lenzen all architectural monuments of the Brandenburg city of Lenzen and its districts are listed.} \\
\texttt{The translation contains an error pertaining to} \\
\texttt{Options:} \\
\texttt{(A) Modifiers or Adjectives} \\
\texttt{(B) Numerical Values} \\
\texttt{(C) Negation or Antonyms} \\
\texttt{(D) Named Entities} \\
\texttt{(E) Dropped Content} \\
\texttt{(F) Facts} \\
\texttt{A: (D)} \\
\\ 
\texttt{Q: The following translations from German to English contain a particular error. That error will be one of the following types: Named Entities: An entity (names, places, locations, etc.) is changed to a different entity. Numerical Values: Numerical values (ordinals or cardinals), dates, and/or units are changed. Modifiers or Adjectives: The modifiers and adjectives pertaining to a noun are changed. Negation or Antonyms: Introduce or remove a negation or change comparatives to their antonyms. Facts: Trivial factual errors not pertaining to the above classes are introduced in the translations. Dropped Content: A significant clause in the translation is removed. Please identify that error.  Source: Auf dieser Seite sind die Baudenkmäler der oberbayerischen Großen Kreisstadt Landsberg am Lech zusammengestellt.} \\
\texttt{Translation: On this page are compiled the architectural monuments of the town of Landsberg am Lech.} \\
\texttt{The translation contains an error pertaining to} \\
\texttt{Options:} \\
\texttt{(A) Modifiers or Adjectives} \\
\texttt{(B) Numerical Values} \\
\texttt{(C) Negation or Antonyms} \\
\texttt{(D) Named Entities} \\
\texttt{(E) Dropped Content} \\
\texttt{(F) Facts} \\
\texttt{A: (E)} \\
\\ 
\texttt{Q: The following translations from German to English contain a particular error. That error will be one of the following types: Named Entities: An entity (names, places, locations, etc.) is changed to a different entity. Numerical Values: Numerical values (ordinals or cardinals), dates, and/or units are changed. Modifiers or Adjectives: The modifiers and adjectives pertaining to a noun are changed. Negation or Antonyms: Introduce or remove a negation or change comparatives to their antonyms. Facts: Trivial factual errors not pertaining to the above classes are introduced in the translations. Dropped Content: A significant clause in the translation is removed. Please identify that error.  Source: Łeba ist eine Kleinstadt und ein Badeort im Powiat Lęborski der polnischen Woiwodschaft Pommern.} \\
\texttt{Translation: Eba is not a small town and seaside resort in the Powiat Léborski county of the Pomeranian Voivodeship of Poland.} \\
\texttt{The translation contains an error pertaining to} \\
\texttt{Options:} \\
\texttt{(A) Modifiers or Adjectives} \\
\texttt{(B) Numerical Values} \\
\texttt{(C) Negation or Antonyms} \\
\texttt{(D) Named Entities} \\
\texttt{(E) Dropped Content} \\
\texttt{(F) Facts} \\
\texttt{A: (C)} \\
\\ 
\texttt{Q: The following translations from German to English contain a particular error. That error will be one of the following types: Named Entities: An entity (names, places, locations, etc.) is changed to a different entity. Numerical Values: Numerical values (ordinals or cardinals), dates, and/or units are changed. Modifiers or Adjectives: The modifiers and adjectives pertaining to a noun are changed. Negation or Antonyms: Introduce or remove a negation or change comparatives to their antonyms. Facts: Trivial factual errors not pertaining to the above classes are introduced in the translations. Dropped Content: A significant clause in the translation is removed. Please identify that error.  Source: Weingarten ist eine Mittelstadt im Südosten Baden{-}Württembergs in Deutschland.} \\
\texttt{Translation: Weingarten is a town in Germany.} \\
\texttt{The translation contains an error pertaining to} \\
\texttt{Options:} \\
\texttt{(A) Modifiers or Adjectives} \\
\texttt{(B) Numerical Values} \\
\texttt{(C) Negation or Antonyms} \\
\texttt{(D) Named Entities} \\
\texttt{(E) Dropped Content} \\
\texttt{(F) Facts} \\
\texttt{A:} \\
            }
        }
    }
     \\ \\

    \theutterance \stepcounter{utterance}  
    & \multicolumn{4}{p{0.6\linewidth}}{
        \cellcolor[rgb]{0.561,0.737,0.561}{
            \makecell[{{p{\linewidth}}}]{
                \texttt{\tiny{[P1$\rangle$GM]}}
                \texttt{To determine the error type, let's compare the source and translation:} \\
\\ 
\texttt{**Source:** "Weingarten ist eine Mittelstadt im Südosten Baden{-}Württembergs in Deutschland."} \\
\texttt{(Translation: Weingarten is a medium{-}sized city in the southeast of Baden{-}Württemberg in Germany.)} \\
\\ 
\texttt{**Translation:** "Weingarten is a town in Germany."} \\
\\ 
\texttt{**Analysis:**} \\
\texttt{1.  **Named Entities:** "Weingarten" and "Germany" are preserved. "Baden{-}Württemberg" is missing, but this is part of the location description.} \\
\texttt{2.  **Numerical Values:** None present.} \\
\texttt{3.  **Modifiers/Adjectives:** "Mittelstadt" (medium{-}sized city) is translated as "town". While slightly less specific, it's not a major modifier error in the context of the other missing info.} \\
\texttt{4.  **Negation/Antonyms:** No negation introduced.} \\
\texttt{5.  **Dropped Content:** The source sentence specifies the location as "im Südosten Baden{-}Württembergs" (in the southeast of Baden{-}Württemberg). The translation completely omits "Baden{-}Württemberg" and the directional "southeast". This is a significant clause describing the specific location within Germany.} \\
\texttt{6.  **Facts:** The remaining facts are correct.} \\
\\ 
\texttt{The most significant error is the omission of the specific region (Baden{-}Württemberg) and its location within that region} \\
            }
        }
    }
    & & \\ \\

    \theutterance \stepcounter{utterance}  
    & & & \multicolumn{2}{p{0.3\linewidth}}{
        \cellcolor[rgb]{0.9,0.9,0.9}{
            \makecell[{{p{\linewidth}}}]{
                \texttt{\tiny{[GM$|$GM]}}
                \texttt{Target: (E)} \\
            }
        }
    }
    & & \\ \\

    \theutterance \stepcounter{utterance}  
    & & & \multicolumn{2}{p{0.3\linewidth}}{
        \cellcolor[rgb]{0.9,0.9,0.9}{
            \makecell[{{p{\linewidth}}}]{
                \texttt{\tiny{[GM$|$GM]}}
                \texttt{game\_result = LOSE} \\
            }
        }
    }
    & & \\ \\

\end{supertabular}
}

\end{document}
